\begin{russian}
Персистентная гомология надёжно выявляет временную рекуррентную структуру в траекториях нейросетевых аудиоэмбеддингов. Однако эта структура лишь слабо соответствует семантике музыкальных жанров, сильно зависит от выбора модели эмбеддингов и локализована на повторяющихся музыкальных секциях лишь для меньшинства треков. В~данной работе строится воспроизводимый пайплайн на датасете GTZAN (999~треков, 10~жанров) с~четырьмя представлениями (MERT, MuQ, EnCodec, MIR) и~тремя суррогатными контролями. Топологические признаки классифицируют жанр на уровне случая (macro-F1~$\approx 0{,}22$), тогда как усреднённые эмбеддинги достигают~$0{,}82$, что демонстрирует разделение между топологией траекторий и~жанровой семантикой. Работа выявляет методологические ловушки, включая почти двукратную переоценку доли <<значимых>> циклов при использовании коцикла вместо гомологического представителя. Вклад работы~--- карта применимости TDA к~аудиоэмбеддингам и~строгий методологический шаблон; отрицательные результаты являются центральным вкладом.
\end{russian}
